% !TeX root = ../../Book2.tex
\section{\texorpdfstring{\(L^p\)}{Lᵖ} Fourier 变换的局限}
\label{b2:sec:0403}

把指数从 \(2\) 往两侧移动，会改变变换所能取值的空间。下面用尺度和相互分离的波包解释为什么 \(p>2\) 不能照搬上一节的结论，并为分布意义下的变换留下明确的问题。

\subsection{\texorpdfstring{\(p>2\)}{p>2} 时的有界延拓障碍}
\label{b2:subsec:040301}

先作尺度检验。假设存在 \(\|\widehat f\|_q\leq C\|f\|_p\) 对全部光滑紧支集函数成立。固定非零 \(\varphi\in C_c^\infty\)，其变换快速下降，故所需的各个范数均有限且非零。对 \(f_a(x)=\varphi(ax)\)、\(a>0\)，伸缩公式给出
\[
 \|f_a\|_p=a^{-d/p}\|\varphi\|_p,
 \quad\|\widehat {f_a}\|_q=a^{-d+d/q}\|\widehat\varphi\|_q,
\]
无穷指数按倒数为零解释。将两式代入假定的估计并消去固定正数，得到
\[
 a^{d(1/p+1/q-1)}\leq C_\varphi\quad(a>0).
\]
若幂为正，让 \(a\to\infty\)；若幂为负，让 \(a\downarrow0\)，均产生矛盾。因此必须有 \(1/p+1/q=1\)。尺度检验只筛去不可能的指数，并不证明剩下的指数可用。

\begin{theorem}\label{b2:thm:04-p-above-two}
设 \(d\) 为正整数，\(2<p\leq\infty\)、\(1/p+1/q=1\)，其中 \(1/\infty=0\)。采用 \(\widehat f(\xi)=\int_{\mathbb R^d}f(x)e^{-2\pi\mathrm{i}x\cdot\xi}\,\mathrm{d}x\) 的归一化，则不存在与 \(f\) 无关的有限常数 \(C\)，使 \(\|\widehat f\|_q\leq C\|f\|_p\) 对所有复值函数 \(f\in C_c^\infty(\mathbb R^d;\mathbb C)\) 成立。
\end{theorem}
\begin{proof}
固定非零 \(\varphi\in C_c^\infty\)，选 \(L\) 足够大，使 \(\varphi(x-jLe_1)\) 的支集互不相交。对符号 \(\epsilon_j\in\{-1,1\}\)，令
\[
 f_N(x)=\sum_{j=1}^N\epsilon_j\varphi(x-jLe_1).
\]
于是 \(\|f_N\|_p=N^{1/p}\|\varphi\|_p\)，\(p=\infty\) 时取 \(N^{1/p}=1\)；其变换为 \(\widehat\varphi(\xi)\sum_{j=1}^N\epsilon_je^{-2\pi\mathrm{i}jL\xi_1}\)。

把全部 \(2^N\) 种符号等权平均。这里 \(\mathbb E\) 只是有限和除以 \(2^N\)，\(\mathbb P\) 只是所占符号组的比例。对任意模为 \(1\) 的复数 \(z_j\)，记 \(Z=\sum_j\epsilon_jz_j\)、\(X=|Z|^2\)。翻转某个符号会使含它奇数次的单项变号，因此这种单项的平均为零。于是
\[
 \mathbb E|Z|^2=\sum_{j,k}z_j\overline{z_k}\,\mathbb E(\epsilon_j\epsilon_k)=N.
\]
展开 \(|Z|^4\) 时，一个四重指标只有每个指标出现偶数次才能保留。它必属于三个成对划分之一：第一位置与第二、第一与第三、或第一与第四配对。每种划分至多有 \(N^2\) 个项，每项系数模为 \(1\)。全相等的指标可能被多次计入，但我们只求上界，故 \(\mathbb EX^2\leq3N^2\)。令 \(E=\{X\geq N/2\}\)，在其补集上 \(X<N/2\)，而在 \(E\) 上使用 Cauchy--Schwarz，得
\[
 N=\mathbb EX\leq N/2+(\mathbb EX^2)^{1/2}\mathbb P(E)^{1/2},
\]
因此 \(\mathbb P(E)\geq1/12\)，进而 \(\mathbb E|Z|^q\geq12^{-1}(N/2)^{q/2}\)。

逐频率取 \(z_j=e^{-2\pi\mathrm{i}jL\xi_1}\)，乘以 \(|\widehat\varphi(\xi)|^q\) 再积分。有限平均可与积分交换，得到
\[
 2^{-N}\sum_{\epsilon\in\{-1,1\}^N}\|\widehat {f_N}\|_q^q
 \geq12^{-1}(N/2)^{q/2}\|\widehat\varphi\|_q^q.
\]
至少一个被平均的非负数不小于平均值，故有一组选定符号使
\[
 \|\widehat {f_N}\|_q\geq c_qN^{1/2}\|\widehat\varphi\|_q.
\]
光滑紧支集函数的变换快速衰减且非零，故右侧范数有限而正。若所求估计成立，则 \(N^{1/2-1/p}\) 应一致有界，与 \(p>2\) 矛盾。
\end{proof}

这个反例甚至使用光滑紧支集函数，故障碍来自范数估计本身，并非某个积分未定义。它否定的是在整个 \(L^p\) 上取值于 \(L^{p'}\) 的有界延拓；并没有否定特殊函数仍可能有很好的变换。

\subsection{普通积分定义的局限}
\label{b2:subsec:040302}

在无限测度空间上，较大的 \(p\) 不意味着 \(L^1\)。例如在 \(\mathbb R^d\) 上，对 \(2<p<\infty\) 取 \(d/p<\alpha\leq d\)，则 \((1+|x|)^{-\alpha}\in L^p\setminus L^1\)；对 \(p=\infty\)，任取 \(0<\alpha\leq d\) 也得到有界但不可积的函数。乘上单位模复指数以后，绝对可积性没有改善，所以定义式不能作为普通 Lebesgue 积分使用。这并不等于某种指定的振荡截断必定发散，而是说普通积分尚未提供定义。

\ProseParagraphBreak

有时可以指定对称截断、Gauss 正则化或其他极限。不同极限的存在性各需证明；一般的 \(L^p\) 等价类也没有预先规定的点值。即使在 \(L^2\) 中，构造的是频率侧的等价类；一句“取每个频率的极限”并不能替代范数完备性给出的定义。

\subsection{\texorpdfstring{\(L^2\)}{L²} Fourier 乘子及其推广障碍}
\label{b2:subsec:040303}

\begin{definition}\label{b2:def:04-multiplier}
设 \(d\) 为正整数。对 \(m\in L^\infty(\mathbb R^d;\mathbb C)\)、\(f\in L^2(\mathbb R^d;\mathbb C)\)，以 \(2\pi\) 归一化的 Plancherel 变换 \(\widehat f=\mathcal F_2f\) 定义
\[
 T_mf\coloneqq\mathcal F_2^{-1}(m\widehat f),
\]
称为符号 \(m\) 的\term[fourier-chengzi-suanzi]{Fourier 乘子算子}[Fourier multiplier operator]。
\end{definition}
\symidx{\(T_m\)：符号为 \(m\) 的 Fourier 乘子算子}

\begin{proposition}\label{b2:prop:04-multiplier-norm}
设 \(d\) 为正整数，\(\mathcal F_2\) 为 \(2\pi\) 归一化的 Plancherel 变换，\(m\in L^\infty(\mathbb R^d;\mathbb C)\)，\(T_mf=\mathcal F_2^{-1}(m\mathcal F_2f)\) 定义于 \(L^2(\mathbb R^d;\mathbb C)\)。有 \(\|T_m\|_{L^2\to L^2}=\|m\|_\infty\)。若 \(K\in L^1(\mathbb R^d;\mathbb C)\)、\(m=\widehat K\) 几乎处处成立，则 \(T_mf=K*f\) 对每个 \(f\in L^2(\mathbb R^d;\mathbb C)\) 成立，等式按 \(L^2\) 等价类理解；卷积公式还定义了每个 \(L^p(\mathbb R^d;\mathbb C)\)、\(1\leq p\leq\infty\) 上的有界算子，范数至多为 \(\|K\|_1\)。
\end{proposition}
\begin{proof}
酉性给出上界。当 \(\|m\|_\infty=0\) 时两边均为零；否则，对任意 \(0<a<\|m\|_\infty\)，集合 \(\{|m|>a\}\) 有正测度。用有限球穷尽空间，可在其中取一个有限正测度子集 \(E\)。由满射性可选 \(f\in L^2\) 使 \(\widehat f=\mathbf1_E\)，于是 \(\|T_mf\|_2>a\|f\|_2\)，令 \(a\uparrow\|m\|_\infty\) 得到等号。卷积表示先对 \(f\in L^1\cap L^2\) 由\cref{b2:thm:03-convolution-transform}成立：Young 不等式保证 \(K*f\) 同时属于 \(L^1\) 与 \(L^2\)，从而可积变换与 Plancherel 变换的一致性可用于两侧。再取 \(f_n\in L^1\cap L^2\) 在 \(L^2\) 中逼近任意 \(f\)；Young 估计使 \(K*f_n\to K*f\)，已证乘子界使 \(T_mf_n\to T_mf\)，故等式推广到 \(L^2\)。其余指数直接由\cref{b2:thm:03-young}给出。
\end{proof}

对 \(p\ne2\)，仅有 \(m\) 有界通常不足以取得同样的乘子结论；符号的平滑程度、奇性与维数会参与估计。\chapref{b2:ch:22}将研究这些性质怎样决定乘子的有界指数范围。


\ProseParagraphBreak

平滑频率截断说明为何符号的导数有用。取 \(m\in C_c^\infty(\mathbb R^d;\mathbb C)\)，令 \(K=\check m\)。\cref{b2:prop:03-higher-derivative-decay}将可积的高阶纯坐标导数转成频率衰减；选取大于维数的阶数，可得 \(K\in L^1\)，所以 \(T_m\) 在每个 \(L^p\) 上都由 Young 不等式控制。对 \(R>0\) 令 \(m_R(\xi)=m(\xi/R)\)，相应核为 \(K_R(x)=R^dK(Rx)\)，且 \(\|K_R\|_1=\|K\|_1\)。于是这一族算子对全部尺度具有同一个范数界。

\begin{proposition}[平滑频率逼近]\label{b2:prop:04-smooth-frequency-approximation}
设 \(d\) 为正整数，\(m\in C_c^\infty(\mathbb R^d;\mathbb C)\)、\(m(0)=1\)，按正号 Fourier 积分定义 \(K=\check m\)，并对每个 \(R>0\) 令 \(K_R(x)=R^dK(Rx)\)。对 \(1\leq p<\infty\)、\(f\in L^p(\mathbb R^d;\mathbb C)\)，当 \(R\to\infty\) 时有 \(K_R*f\to f\) 于 \(L^p\)；若复值函数 \(f\) 在 \(\mathbb R^d\) 上有界且一致连续，则收敛为一致收敛。对 \(f\in L^2(\mathbb R^d;\mathbb C)\)，这个卷积等于 \(T_{m(\cdot/R)}f\)。
\end{proposition}
\begin{proof}
先核对核的可积性。取整数 \(N>d\)，则 \(m\) 及各 \(\partial_j^Nm\) 都属于 \(L^1\)；由\cref{b2:prop:03-higher-derivative-decay}以及 \(\check m(x)=\widehat m(-x)\)，有 \(|K(x)|\leq C|x|^{-N}\) 对 \(|x|\geq1\) 成立，而在单位球内 \(|K|\leq\|m\|_1\)。两部分合起来给出 \(K\in L^1\)。对光滑紧支集的 \(m\) 使用上一章的反演，再作反射换元，得到 \(\widehat K=m\)；这里两边连续，所以等式处处成立。特别 \(\int K=\widehat K(0)=m(0)=1\)。注意 \(K\) 可以有复值且不必非负，所需的是绝对值的统一质量界与尾部集中性。将卷积换元，得到
\[
 K_R*f-f=\int K(y)\bigl(f(\cdot-y/R)-f\bigr)\,\mathrm{d}y.
\]
因而由积分型 Minkowski 不等式，
\[
 \|K_R*f-f\|_p
 \leq\int|K(y)|\,\|f(\cdot-y/R)-f\|_p\,\mathrm{d}y.
\]
对每个固定 \(y\)，有限指数的平移连续性使内层范数趋于零；它又不超过 \(2\|f\|_p\)。控制收敛给出结论。若 \(f\) 有界一致连续，用同一论证把 \(L^p\) 范数换成上确界范数即可。\(L^2\) 上的乘子识别来自前一命题，因为 \(\widehat {K_R}(\xi)=m(\xi/R)\)。
\end{proof}

这个证明也显示为何不能把结论无条件推广到 \(L^\infty\) 范数。在一维取阶跃函数 \(f=\mathbf1_{(0,\infty)}\)，它与光滑核卷积后具有连续代表。若一个连续函数与 \(f\) 的本性上确界距离小于 \(1/2\)，可取统一的 \(c<1/2\)，使它在左侧几乎处处距 \(0\) 至多为 \(c\)，在右侧几乎处处距 \(1\) 至多为 \(c\)。连续性使这些闭圆盘约束在原点也同时成立，但两圆盘不相交。因此 \(\|K_R*f-f\|_\infty\geq1/2\)；这里的收敛限制由函数的跳跃造成，尽管每个算子的 \(L^\infty\) 范数都统一有界。

\ProseParagraphBreak

相反，对 \(R>0\)，一维的硬截断符号 \(m=\mathbf1_{[-R,R]}\) 对应
\[
 K(x)=\dfrac{\sin(2\pi Rx)}{\pi x},
\]
该核不属于 \(L^1\)。对 \(f\in C_c^\infty\)，仍可用有限频率积分和 Fubini 验证 \(T_mf=K*f\)，但不能再用 \(\|K\|_1\) 给出有限的 Young 上界。\(L^2\) 上的范数仍为 \(1\)，因为符号是有界的正交投影。这个例子把两种方法的职责分开：酉性直接控制平方能量，绝对可积核则提供全部指数上的一种充分条件。

\subsection{分布意义的 Fourier 变换的动机}
\label{b2:subsec:040304}

若不能要求变换是一个普通函数，还可以要求它对足够好的测试函数产生确定的数值。对 \(f\in L^p\) 和光滑紧支集的 \(\varphi\)，函数 \(\widehat\varphi\) 快速衰减，因而 \(\int f\widehat\varphi\) 总有意义。若 \(f\) 原来可积，Fubini 又给出
\[
 \int\widehat f(\xi)\varphi(\xi)\,\mathrm{d}\xi
 =\int f(x)\widehat\varphi(x)\,\mathrm{d}x.
\]
右侧因此提供了延拓变换的线索，且对每个 \(1\leq p\leq\infty\) 都适用。下一章将选取在 Fourier 变换下保持不变的 Schwartz 测试空间，使这个配对成为完整定义。

\ProseParagraphBreak

测试函数的下降速度在这里承担了可积性责任。例如取整数 \(N\) 满足 \(Np'>d\)，当 \(p'=\infty\) 时只需取任意正整数。则 Hölder 不等式给出
\[
 \left|\int f\widehat\varphi\right|
 \leq C_{N,p,d}\|f\|_p
 \sup_x(1+|x|)^N|\widehat\varphi(x)|.
\]
其中 \(p'=1\) 也由 \(N>d\) 包含在内。反复分部积分把右侧的加权上确界控制在 \(\varphi\) 的有限多个导数与紧支集大小之内。这样一个可能没有可积频率代表的对象，仍能对每个测试函数给出稳定的线性数值；下一章要精确说明的是这种稳定性应在哪个测试函数拓扑中表述。
